\section{Conclusion}
This paper introduced LARA, a framework that leverages Linguistic-Adaptive Retrieval-Augmentation to address multi-turn intent classification challenges through zero-shot settings across multiple languages. Unlike other supervised Fine-Tuning (SFT) models, which require a hard-to-collect multi-turn dialogue set, our method requires only a single-turn training set to train a conventional model, combining it with an innovative in-context retrieval augmentation for multi-turn intent classification. LARA substantially improved user satisfaction by 1\% over multi-turn sessions, reduced the transfer-to-agent rate by 0.5\% and saved the cost of hundreds of agents, which is a key business metric in our industrial application.

% marking a significant advancement in the field of conversational AI.

The empirical results underscore LARA's capability to enhance intent classification accuracy by 3.67\% over existing methods while reducing inference time, thus facilitating real-time application adaptability. Its strategic approach to managing extensive intent varieties without exhaustive dataset requirements presents a scalable solution for complex, multi-lingual conversational systems.

For future work, we intend to extend and apply the LARA framework to recommendation tasks in other domains, such as understanding how user intent may shift during a POI or itinerary recommendation for tourism purposes~\cite{halder2024survey}. This will enable us to better capture evolving user preferences due to temporal shifts, changing contexts, and individual or group behavioural patterns, which is also applicable to general sequence recommendations.

\vspace{3mm}
{
{\noindent\bf Acknowledgments.} 
This research is supported in part by the Ministry of Education, Singapore, under its Academic Research Fund Tier 2 (Award No. MOE-T2EP20123-0015). Any opinions, findings and conclusions, or recommendations expressed in this material are those of the authors and do not reflect the views of the Ministry of Education, Singapore.
}